% ============================================================
%  परिशिष्ट ग : सन्दर्भ एवं अग्रिम पठन (References & Further Reading)
% ============================================================
\chapter*{परिशिष्ट ग: सन्दर्भ एवं अग्रिम पठन}
\addcontentsline{toc}{chapter}{परिशिष्ट ग: सन्दर्भ एवं अग्रिम पठन}
\markboth{सन्दर्भ एवं अग्रिम पठन}{सन्दर्भ एवं अग्रिम पठन}

\noindent नीचे विश्वसनीय एवं निःशुल्क स्रोतों की सूची दी गई है, जिनसे विद्यार्थी इस
पुस्तक के विषयों को और गहराई से समझ सकते हैं। (सभी लिंक पर क्लिक किया जा सकता है।)

\section*{सामान्य परिचय (General Introduction to AI)}
\begin{itemize}[leftmargin=1.4em]
  \item \href{https://www.elementsofai.com/}{Elements of AI} -- निःशुल्क ऑनलाइन
        पाठ्यक्रम (free online course), University of Helsinki.
  \item \href{https://www.youtube.com/watch?v=aircAruvnKK}{3Blue1Brown: Neural
        Networks} -- तंत्रिका जाल का दृश्य परिचय (visual introduction).
  \item \href{https://ai.google/education/}{Google AI Education} -- AI एवं ML
        सीखने हेतु सामग्री।
\end{itemize}

\section*{निःशुल्क उपकरण (Free Hands-on Tools)}
\begin{itemize}[leftmargin=1.4em]
  \item \href{https://teachablemachine.withgoogle.com/}{Teachable Machine} --
        बिना कोड के image/sound classification.
  \item \href{https://colab.research.google.com/}{Google Colab} -- ब्राउज़र में
        निःशुल्क Python.
  \item \href{https://phyphox.org/}{Phyphox} -- मोबाइल सेंसर से भौतिक डेटा।
  \item \href{https://orangedatamining.com/}{Orange Data Mining} -- दृश्य रूप से
        ML समझना।
  \item \href{https://molview.org/}{MolView} -- अणुओं की 3D संरचना।
\end{itemize}

\section*{विषयवार स्रोत (Subject-wise Sources)}
\begin{itemize}[leftmargin=1.4em]
  \item \textbf{भौतिकी:} \href{https://home.cern/}{CERN} एवं
        \href{https://www.nasa.gov/}{NASA} -- कण भौतिकी एवं खगोल में डेटा-विज्ञान।
  \item \textbf{रसायन:} \href{https://pubchem.ncbi.nlm.nih.gov/}{PubChem} --
        रासायनिक यौगिकों का खुला डेटाबेस।
  \item \textbf{जीव विज्ञान:} \href{https://alphafold.ebi.ac.uk/}{AlphaFold
        Protein Structure Database} एवं
        \href{https://www.ncbi.nlm.nih.gov/}{NCBI} -- प्रोटीन एवं जीनोम डेटा।
  \item \textbf{डेटा-विज्ञान:} \href{https://numpy.org/}{NumPy} एवं
        \href{https://matplotlib.org/}{Matplotlib} -- Python पुस्तकालय।
\end{itemize}

\begin{padhe}
इन स्रोतों का उपयोग करते समय \emph{स्रोत की विश्वसनीयता} (credibility) अवश्य जाँचें।
सरकारी, विश्वविद्यालयी एवं प्रतिष्ठित वैज्ञानिक संस्थाओं की वेबसाइट सामान्यतः अधिक
विश्वसनीय होती हैं। AI से प्राप्त उत्तर को भी किसी विश्वसनीय स्रोत से सत्यापित करें।
\end{padhe}

\section*{चित्र-आभार (Image Credits)}
\noindent इस पुस्तक में प्रयुक्त छायाचित्र \emph{सार्वजनिक डोमेन} (Public Domain) से लिए
गए हैं तथा Wikimedia Commons के सौजन्य से उपलब्ध हैं:
\begin{itemize}[leftmargin=1.4em]
  \item \textbf{चित्र~\ref{fig:hubble}} (हबल अल्ट्रा डीप फ़ील्ड): NASA \& ESA --
        Public Domain;
        \href{https://commons.wikimedia.org/wiki/File:Hubble_ultra_deep_field.jpg}{Wikimedia Commons}.
  \item \textbf{चित्र~\ref{fig:caffeine}} (कैफ़ीन अणु): Public Domain;
        \href{https://commons.wikimedia.org/wiki/File:Caffeine-3D-balls.png}{Wikimedia Commons}.
  \item \textbf{चित्र~\ref{fig:aspirin}} (एस्पिरिन अणु): Benjah-bmm27 -- Public
        Domain; \href{https://commons.wikimedia.org/wiki/File:Aspirin-3D-balls.png}{Wikimedia Commons}.
  \item \textbf{चित्र~\ref{fig:dna}} (DNA द्विकुंडलिनी): Forluvoft -- Public
        Domain; \href{https://commons.wikimedia.org/wiki/File:DNA_simple2.svg}{Wikimedia Commons}.
\end{itemize}
\noindent अन्य सभी आरेख (diagrams) लेखक द्वारा \emph{TikZ} में मूल रूप से बनाए गए हैं।
